\documentclass[titlepage,a4paper]{jsarticle}
\usepackage{../../../sty/import}% 各種パッケージインポート
\usepackage{../../../sty/title}% タイトルページの変更

%% タイトルページの変数
% レポートタイトル
\title{心理学特論第5回目出席票}
% 提出日
\expdate{\today}
% 科目名
\subject{心理学特論}
% 分野
\class{情報経営システム工学分野}
% 学年
\grade{M1}
% 学籍番号
\mynumber{24336488}
% 記述者
\author{本間三暉}

\begin{document}
% titleページ作成
\maketitle
\section*{keyword\#1}
認知的ケチ
\section*{keyword\#2}
ヒューリスティック

\section{心理学特論の第05回出席課題は身近な行動経済学としての「心の会計」を考えましょう．}
\subsection{「あぶく銭効果」を実感した自分(または周囲)の経験があれば，経緯を踏まえて書いてください．
経験がない場合は，臨時収入があったとき，どのようにしているかを簡潔に書いてください}
これはValorantというオンラインFPSゲームの話であるが，普段一緒にしないゲームが強い友人とやり，あっさり勝ってしまったときに，
その揺り戻しでそれ以降の試合で相手が強くなっているにも関わらず，行けるだろうみたいな軽いノリでランクマッチに潜ってしまい，気がついたらランクポイントが溶けているという経験がある．

\subsection{「サンクコストの呪縛」は強烈です．これ以上の利得はないだろうと薄々わかっていながら，
「損切り」ができないでずるずると引きずってしまった経験(自分だけでなく，他人のでも可)や，世の中で起きた事象について，経緯を踏まえて簡潔に書いてください．}
  これもまた，ValorantというオンラインFPSゲームの話であるが，あとちょっとでランクが上がるとき，本当に調子が悪く負けてしまい，
ランクが上がるのが遠ざかってしまったときに，今日は調子が悪いからやめようか迷っているが，あとちょっとでランクが上がるからもう少しやろうと考えてランクマッチに潜ってしまい，
結局負けに負けてランクが上がるのが更に遠ざかってしまったという経験が何度もある．

\section{心理学特論課題サイコロワーク2 260507 解答記入欄}
(01)～(20)に該当する適切な用語を解答群から選んで下さい．

% 人間が選択をする時，伝統的な経済学ではたくさんの情報をきちんと吟味してから決断するとされています．
% これを (01) といいます．しかしすべての場面で，そのような判断を行うのは手間や時間がかかりすぎます．
% そのため安いものやこだわりのないものを購入する場合や，瞬時に判断が必要となる場合，情報が大きすぎる場合は効率良く決断をする (02) を対応します．
% (02) は常に適切な判断に結びつくとは限りません．
% 見たいものだけを見て，聞きたいものだけを聞くといった，都合よく解釈するなど誤った考え方，すなわち (03) を引き起こすこともあります．
% どうして興味のある事だけ聞こえるのかというと，人は選択的に知覚できるからです．
% 例えば，人混みでも自分の関心のある言葉は聞き取れるという (04) も，(05) の (02) として分類できます．

% ＝ (01)～(05) 解答群 ＝
% ＊問題に直接関係のない言葉も含まれています．
% システマティック
% オートマティック
% ヒューリスティック
% バイアス
% ノイエス
% カクテルパーティー効果
% フレーミング効果
% 代表性
% 係留性
% 利用可能性

% お金の使いかたはどのように決めているのでしょうか．実は手に入れた方法で遣い方が変わります．
% 損得や確率が関係する (06) にはその (06) のために状況を評価する (07) という段階が影響を与えています．
% ここでの検討項目の1つに心の口座があります．
% 頭の中で，生活費，娯楽費のような複数の口座を作り，それぞれの口座の中で (08) を最大にしようと，無意識のうちにやりくりをしているという考え方です．
% これをメンタル・アカウンティング，すなわち (09) といいます．どのようにして得たお金かによっても使いかたは変わります．
% 合理的に考えればどのお金も意味は同じはずですが，実際には働いて稼いだお金は大切に使う傾向があるのに対して，ギャンブルなどで苦労せず手に入れたお金の使いかたは無造作になりがちです．
% このような傾向を (10) といいます．人は手元のお金に意味を付け加えそれをもとに使いかたを決めるのです．

% ＝ (06)～(10) 解答群 ＝
% ＊問題に直接関係のない言葉も含まれています．
% 傍観者効果
% ハウスマネー効果
% ヴェブレン効果
% 費用対効果
% 意思決定
% 心の会計
% 粉飾決算
% 後処理
% 前処理
% 抽選確率
% 判断の誤謬


% 有名人など高い評価を得ている人物と自分とを関連付けたくなることを心理学では (11) という．
% 「すごい人と関係のある自分もすごい」と周囲にアピールして，自己評価を上げたいという心理作用である．
% 自分の母校や出身地，好きな野球チームを自慢するのも一緒の (11) だが，世間的に評価が低い友人・知人とはそんなに親しくないよとアピールして距離を取る．
% 同期の同僚が出世した場合，自分が同期より組織から評価されていないという見方もできる．つまり同期の出世によって自己評価が下がったと感じてストレスが高まる．
% テッサーは，こうした自己評価の心理を『自己評価維持モデル』として理論化した．
% 『自己評価維持モデル』では，相手との (12) ，(13) ，(14) の3要素で自分を評価すると考える．(13) は，課題との関与度で，(14) は，いわば成績である．
% 同窓生が有名人の場合，(12) が近く，相手の (14) が高く，(13) は低いから (11) がはたらく．
% 同期の同僚も (12) は近いし，相手の (14) も高いのだが，組織での役割分担，すなわち会社の仕事や研究テーマといった同じような課題に取り組むために (13) が高くなる．
% このような場合はストレスを感じてしまう．
% 自己評価が下がった場合の調整法としては，(12) については相手を遠ざけるなどして距離を置く，(14) については自分が努力して成績を上げる，または相手の足を引っ張るなどして貶める，
% (13) については課題への関心・執着を低くするなどが考えられる．
% とは言え，人は無意識のうちに自分と他人とを比較し，自分の能力を正確に評価しようとしている．
% 心理学ではこれを (15) といい，自分より能力が上の人と比較することを上方比較，能力が下の人と比較することを下方比較という．

% ＝ (11)～(15) 解答群 ＝
% ＊問題に直接関係のない言葉も含まれています．
% 物理的距離
% 心理的距離
% 自己関連性
% 他者関連性
% 後光効果
% 栄光浴
% 文化的比較理論
% 社会的比較理論
% 遂行
% 同行

% 何かで思ったようにうまく行かなかったとき，その原因をどこに求めますか？「才能」，「努力」，「運」，「難易度」など様々な要因が考えられます．このように，行動の結果として原因を何に求めるかを「(16)」といいます．
% (17) は「成功や失敗そのもの」よりも，「その原因をどのように捉えるか」が重要と述べています．
% 大きなことを成し遂げられる人は，原因を「自分で変えることができること」すなわち自分の努力に求める傾向が強いです．
% しかし，その反面，自分のせいにしてしまうことで落ち込みやすい面があるのも事実です．
% ちょっとしたことでショックを受けて，傷ついた気持ちを長く引きずってしまう人は「(18)」が低い傾向にあるとも言えます．
% 「(19)」では，(18) の高い人は自分を良く評価してくれる人を好み，(18) の低い人は自分をあまり評価しない人を好む傾向があるとしています．
% つまり，褒め言葉を素直に喜べるかどうかは (18) の一種のバロメーターです．

% 「自信がある人ほど反省しない」面もあるそうです．
% 自分を責め後悔するだけの反省では行動力の低下を招くだけです．
% ベイアスらの調査では，(20) の大半が「周囲に相談しない」で行動しているそうです．
% 何かをするかしないかで他人に相談する時は迷っていて，「やりたくない」「止めて欲しい」という気持ちが隠れていることもあります．

% ＝ (16)～(20) 解答群 ＝
% ＊問題に直接関係のない言葉も含まれています．
% メラマード
% ワイナー
% 行動力のない人
% 行動力のある人
% 原因帰属
% 自己評価
% 客観評価
% 結果帰属
% 認知的不協和理論
% 認知的斉合性理論


\begin{enumerate}
  %% 1~5
  \item システマティック
  \item ヒューリスティック
  \item バイアス
  \item カクテルパーティー効果
  \item 利用可能性

  %% 6~10
  \item 意思決定
  \item 前処理
  \item 費用対効果
  \item 心の会計
  \item ハウスマネー効果
  
  %% 11~15
  \item 栄光浴
  \item 心理的距離
  \item 自己関連性
  \item 遂行
  \item 社会的比較理論
  
  %% 16~20
  \item 原因帰属
  \item ワイナー
  \item 自己評価
  \item 認知的斉合性理論
  \item 行動力のある人
\end{enumerate}

\end{document}

% 